\subsection{Procedimiento reproducible de toma de datos}
\label{sec:protocolo_a}
El procedimiento mantiene unidos el caso previsto y lo que efectivamente sucedió. La revisión de cada ejecución sigue este orden.
\begin{enumerate}
\item Registrar condición, repetición, entradas, intermedios y finales. Identificar qué contraste usa el caso y qué controles deben coincidir con su base.
\item Revisar versiones, recursos, pool y cola. Corroborar qué máquinas están disponibles y qué requisitos se declararon. Anotar la posición del envío dentro de la secuencia.
\item Revisar el plan antes del envío. Comprobar qué tareas pertenecen a cada job y qué dependencias, archivos y finales declara. La agrupación se confirma en las descripciones efectivas.
\item Ejecutar y preservar el identificador, las marcas de inicio y fin de DAGMan, la ubicación e intento de cada job, y los contratos y registros de entrada y salida.
\item Reconstruir las dos cuentas. Para cada obligación, identificar el archivo, su tamaño y el destino. Para cada entrega, vincular archivo y dirección con el intento y conciliar conteos y bytes. Registrar los conflictos de cobertura o identidad.
\item Resolver elegibilidad y comparar con la base. Conservar por separado los casos que incumplen el perfil y sus diagnósticos. Presentar ambos volúmenes, la razón exacta y la duración del mismo workflow.
\end{enumerate}

La duración del workflow comienza y termina según las marcas de DAGMan. Las tareas de preparación y entrega que ocurren dentro de ese intervalo forman parte de la duración. La exportación del historial y la invocación de \Analyzer{} posteriores al término se cronometran por separado. Los registros de transferencias corroboran las entregas y sus volúmenes. Su resolución temporal no se utiliza aquí para estimar una tasa de red a partir de intervalos breves.

\begin{table}[htbp]\centering\small
\caption{Reglas de elegibilidad y comprobaciones del análisis}
\begin{tabularx}{\linewidth}{@{}P{3.3cm}YY@{}}
\toprule Comprobación & Criterio & Consecuencia si falla\\\midrule
Condición y controles & Coinciden grafo, tamaños, grupos, ubicación y política con el nivel documentado. & El caso queda fuera de la comparación del perfil fijo.\\
Terminación e intento & El workflow termina y el intento efectivo puede identificarse para cada job. & Se conserva el diagnóstico de ejecución y se revisa la elegibilidad.\\
Cobertura de archivos & Se reconstruyen los archivos necesarios y se respaldan las entregas con evidencia. & El volumen afectado y DME pueden quedar indeterminados.\\
Coincidencia del cálculo & Se comparan las cuentas manuales, el núcleo y el auditor local sobre los mismos insumos; se distingue este cotejo de la elegibilidad de la ejecución. & Se localiza la discrepancia antes de interpretar un contraste.\\
Uso de la duración & Se conserva el tiempo de cada ejecución y la organización efectiva de jobs. & Se acota la conclusión a la evidencia disponible.\\
\bottomrule
\end{tabularx}
\end{table}

De los 74 envíos del bloque A, dos quedaron fuera del perfil fijo. L01-r01 terminó y fue auditado, pero se ajustaron los intervalos del planificador y del negociador mientras el flujo seguía en curso. Sus repeticiones elegibles son r02, r03 y r04. L05-r01 se interrumpió al detener la sesión y quedó sin aceptación. Los dos registros se conservan. Las restantes 72 ejecuciones cumplieron las reglas de inclusión.
